% ============================================================
% SQL Patterns (PostgreSQL) -- master wrapper for sections/30-code/31-sql/.
% Mirrors the sections/30-code.tex <-> sections/30-code/ pattern one
% level down. \input paths are relative to the MAIN document being
% compiled, so paths below are given in full from the project root.
% ============================================================
\subsection{SQL Patterns (PostgreSQL)}
% ============================================================

Reference for recurring query shapes. Each entry states the problem phrasing
that signals it and the mechanism that makes it work.

\begin{keybox}[Reading a problem]
Most prompts map to a small set of shapes, and the wording is the tell:
\begin{itemize}[label=$\circ$, itemsep=1ex]
  \item ``per group'', ``for each'' $\to$ \texttt{GROUP BY}, or a
  \texttt{PARTITION BY} window when the whole row must survive.
  \item ``top-$N$ per group'', ``most recent per'' $\to$ rank-then-filter,
  or \texttt{DISTINCT ON} when $N = 1$.
  \item ``change from the previous'', ``period-over-period'', ``return''
  $\to$ \texttt{LAG} / \texttt{LEAD}.
  \item ``running'', ``cumulative'', ``rolling'', ``moving average'' $\to$
  an ordinary aggregate with \texttt{OVER}.
  \item ``consecutive'', ``streak'', ``gap'' $\to$ gaps and islands.
  \item ``one column per category'', ``pivot'' $\to$ conditional aggregation.
\end{itemize}

\vspace{1ex}
The first decision on any prompt is whether the answer collapses rows
(\texttt{GROUP BY}) or keeps every row and appends a computed column (a
window function).
\end{keybox}

\begin{keybox}[Logical processing order]
A query is evaluated in this order, regardless of written order:
\[
  \texttt{FROM/JOIN} \to \texttt{WHERE} \to \texttt{GROUP BY} \to
  \texttt{HAVING} \to \texttt{SELECT} \to \texttt{DISTINCT} \to
  \texttt{ORDER BY} \to \texttt{LIMIT}
\]
Two consequences drive most of the patterns below. \texttt{WHERE} runs
before aggregation, so it cannot see aggregates or \texttt{SELECT} aliases.
Window functions are computed at the \texttt{SELECT} step, so they also
cannot appear in \texttt{WHERE} and must be wrapped in a subquery or CTE
to filter on.
\end{keybox}

\subsubsection{Filtering and Joins}

The foundation. Every later pattern assumes you can already restrict rows and
combine tables.

\paragraph{WHERE vs HAVING.} \texttt{WHERE} filters individual rows before
grouping. \texttt{HAVING} filters whole groups after aggregation and is the
only place an aggregate predicate is legal. A filter on a plain
(non-aggregated) column belongs in \texttt{WHERE} so rows are dropped before
the grouping work happens. Reserve \texttt{HAVING} for conditions like
\texttt{COUNT(*) > 5}.

\begin{lstlisting}[style=sql, caption={WHERE prunes rows, HAVING prunes groups}]
SELECT department, COUNT(*) AS headcount
FROM employees
WHERE active = true          -- row filter, before grouping
GROUP BY department
HAVING COUNT(*) > 5;         -- group filter, on an aggregate
\end{lstlisting}

\vspace{-1ex}

\newpage

\paragraph{Joins.} The recognition cue is which side you want to keep when a
match is missing.
\begin{itemize}[label=$\circ$]
  \item \textbf{INNER:} rows with a match on both sides.
  \item \textbf{LEFT / RIGHT:} keep all rows from one side, \texttt{NULL}-fill the other.
  \item \textbf{FULL OUTER:} keep unmatched rows from both sides.
  \item \textbf{CROSS:} Cartesian product, every left row against every right row.
  \item \textbf{SELF:} a table joined to itself under an alias (employee to manager, row to prior row).
  \item \textbf{Semi-join} (\texttt{EXISTS}, \texttt{IN}): keep left rows with at least one match. No row duplication, no right-side columns.
  \item \textbf{Anti-join} (\texttt{NOT EXISTS}, or \texttt{LEFT JOIN ... WHERE right.key IS NULL}): keep left rows with no match.
\end{itemize}

Semi- and anti-joins filter the left table on whether a match exists, without
pulling right-side columns or duplicating rows. Prefer \texttt{NOT EXISTS} over
\texttt{NOT IN} for the anti-join, since a \texttt{NULL} on the right poisons
\texttt{NOT IN} to zero rows.

\begin{lstlisting}[style=sql, caption={Semi-join keeps matches, anti-join keeps non-matches}]
-- Semi-join: customers with at least one order
SELECT c.customer_id, c.name
FROM customers c
WHERE EXISTS (SELECT 1 FROM orders o
              WHERE o.customer_id = c.customer_id);

-- Anti-join: customers with no orders
SELECT c.customer_id, c.name
FROM customers c
WHERE NOT EXISTS (SELECT 1 FROM orders o
                  WHERE o.customer_id = c.customer_id);

-- Anti-join, LEFT JOIN form: unmatched rows are the NULL-filled ones
SELECT c.customer_id, c.name
FROM customers c
LEFT JOIN orders o ON o.customer_id = c.customer_id
WHERE o.customer_id IS NULL;
\end{lstlisting}
\vspace{-4ex}

\subsubsection{Aggregation}

\paragraph{Conditional aggregation (pivot).} The signal is a request for one
column per category on a single row. Fold a categorical column into counts by
wrapping a \texttt{CASE} in an aggregate. Rows the \texttt{CASE} does not match
contribute 0 to \texttt{SUM}. In Postgres the \texttt{FILTER} clause does the
same thing and reads more cleanly.

\begin{lstlisting}[style=sql, caption={Pivot risk categories to columns}]
SELECT
    inspection_type,
    COUNT(*) FILTER (WHERE risk_category = 'Low Risk')      AS low_risk,
    COUNT(*) FILTER (WHERE risk_category = 'High Risk')     AS high_risk,
    COUNT(*)                                                AS total
FROM sf_restaurant_health_violations
GROUP BY inspection_type
ORDER BY total DESC;
\end{lstlisting}

\vspace{-1ex}

The \texttt{SUM(CASE WHEN ... THEN 1 ELSE 0 END)} form is equivalent and
portable across engines. \texttt{COUNT} ignores \texttt{NULL}, so
\texttt{COUNT(CASE WHEN cond THEN 1 END)} (no \texttt{ELSE}) also counts only
the matches.

\subsubsection{NULL Handling}

\texttt{NULL} means unknown, so it propagates through arithmetic and poisons
comparisons. The two functions below are the standard guards.

\paragraph{COALESCE and NULLIF.} \texttt{COALESCE(a, b, c, ...)} returns the
first non-\texttt{NULL} argument, evaluated left to right, or \texttt{NULL} if
all are \texttt{NULL}. The common use is a fallback that keeps a \texttt{NULL}
from poisoning arithmetic: \texttt{salary + commission} is \texttt{NULL}
whenever \texttt{commission} is, so write \texttt{salary + COALESCE(commission, 0)}.
It also supplies a default down a priority list, as in
\texttt{COALESCE(nickname, first\_name, 'Unknown')}. It is ANSI standard,
unlike the vendor-specific \texttt{IFNULL} (MySQL) and \texttt{ISNULL}
(SQL Server), and it accepts any number of arguments.

\texttt{NULLIF(a, b)} is the mirror image: it returns \texttt{NULL} when
\texttt{a = b}, otherwise \texttt{a}. Its standard use guards division,
\texttt{x / NULLIF(y, 0)} yields \texttt{NULL} instead of raising a
divide-by-zero when \texttt{y} is 0.

\begin{lstlisting}[style=sql, caption={NULLIF guards division, COALESCE supplies defaults}]
SELECT
    order_id,
    revenue / NULLIF(quantity, 0)  AS unit_price,  -- NULL, not an error, when quantity = 0
    COALESCE(discount, 0)          AS discount      -- NULL discount treated as 0
FROM orders;
\end{lstlisting}

\vspace{-1ex}

A subtlety with aggregates: \texttt{SUM}, \texttt{AVG}, and
\texttt{COUNT(col)} skip \texttt{NULL}s. So \texttt{AVG(commission)} averages
only the non-\texttt{NULL} rows, while \texttt{AVG(COALESCE(commission, 0))}
counts the \texttt{NULL}s as zeros and enlarges the denominator. 

\paragraph{Gotchas.}
\begin{itemize}
  \item \textbf{\texttt{NOT IN} with NULLs:} if the subquery returns any
  \texttt{NULL}, \texttt{x NOT IN (...)} is never true and the query returns
  zero rows. \texttt{x NOT IN (1, NULL)} expands to
  \texttt{NOT (x=1 OR x=NULL)}, and the \texttt{NULL} comparison poisons the
  result to unknown. Use \texttt{NOT EXISTS}, which is null-safe.
  \item \textbf{Integer division truncates:} \texttt{5 / 2} yields \texttt{2}
  because the fractional part is discarded. Cast one operand:
  \texttt{5::numeric / 2} or \texttt{5.0 / 2}. Bites hardest in ratio and
  percentage columns.
  \item \textbf{\texttt{COUNT(col)} vs \texttt{COUNT(*)}:} the former skips
  \texttt{NULL}s in \texttt{col}, the latter counts every row. A common source
  of off-by-$N$ errors.
\end{itemize}

\subsubsection{Window Functions}

The signal for a window function is a per-row result that still depends on
other rows, such as a rank, a neighbor value, or a running total. Where
\texttt{GROUP BY} would collapse the group to one row, a window function keeps
every row and appends the computed column.

Every window function shares the same shape:
\texttt{f() OVER (PARTITION BY ... ORDER BY ...)}. \texttt{PARTITION BY} splits
the rows into independent groups, \texttt{ORDER BY} sets the order within each
group, and the function is evaluated per row. All of them are computed at the
\texttt{SELECT} step, so to filter on their output you wrap them in a subquery
or CTE.

\paragraph{Rank-then-filter with ROW\_NUMBER().} Use this for top-$N$ per
group when you need to keep the full row, since \texttt{GROUP BY} would collapse
the other columns away. Number rows inside each partition, then filter on that
number in an outer query.

Choose the ranker by tie policy: \texttt{ROW\_NUMBER} assigns distinct numbers
and breaks ties arbitrarily, \texttt{RANK} repeats a rank and leaves gaps,
\texttt{DENSE\_RANK} repeats with no gaps.

\begin{lstlisting}[style=sql, caption={Highest-paid employee per department}]
SELECT department, name, salary
FROM (
    SELECT department, name, salary,
           ROW_NUMBER() OVER (PARTITION BY department
                              ORDER BY salary DESC) AS rn
    FROM employees
) ranked
WHERE rn = 1;
\end{lstlisting}

\vspace{-4ex}

\paragraph{Offset access with LAG() and LEAD().} The signal is a comparison
against the previous or next row: a delta, a period-over-period change, or a
return. These read a value from another row relative to the current one,
without a self-join. \texttt{LAG(col, k, default)} returns the value \texttt{k}
rows earlier in the ordered partition, \texttt{LEAD} looks \texttt{k} rows
ahead. Both default to \texttt{k = 1} and return \texttt{NULL} past the
partition edge unless a fill \texttt{default} is supplied. The \texttt{ORDER BY}
inside \texttt{OVER} defines what ``previous'' means and is mandatory.
\texttt{PARTITION BY} resets the sequence at each group boundary, so a lag never
bleeds from one partition into the next.

\begin{lstlisting}[style=sql, caption={Daily returns per ticker with LAG}]
SELECT ticker, trade_date, close,
       close / LAG(close) OVER (PARTITION BY ticker
                                ORDER BY trade_date) - 1 AS daily_return
FROM prices;
\end{lstlisting}

\vspace{-1ex}

The first row of each partition has no predecessor, so its \texttt{LAG} and the
derived return are \texttt{NULL}, which is correct since there is no prior price
to compare against. \texttt{LEAD} is the same tool aimed forward, used below for
gap detection.

\paragraph{Running and rolling aggregates.} The signal is ``running'',
``cumulative'', or ``moving average''. An ordinary aggregate gains a window when
you attach \texttt{OVER (... ORDER BY ...)}. With an \texttt{ORDER BY} and no
explicit frame, the frame defaults to everything from the partition start up to
the current row, giving a running total. Add an explicit \texttt{ROWS BETWEEN}
frame to get a fixed moving window instead.

\begin{lstlisting}[style=sql, caption={20-day moving average and running volume}]
SELECT ticker, trade_date, close,
       AVG(close) OVER (PARTITION BY ticker ORDER BY trade_date
                        ROWS BETWEEN 19 PRECEDING AND CURRENT ROW) AS ma_20,
       SUM(volume) OVER (PARTITION BY ticker ORDER BY trade_date)  AS cum_volume
FROM prices;
\end{lstlisting}

\vspace{-1ex}

\texttt{ma\_20} is the current row plus the 19 before it, and \texttt{cum\_volume}
uses the default frame for a running total. These are the SQL forms of a pandas
\texttt{.rolling(20).mean()} and \texttt{.expanding().sum()}.

One subtlety: the default frame is \texttt{RANGE UNBOUNDED PRECEDING}, which
includes every row tied on the \texttt{ORDER BY} key, not just rows physically
up to the current one. When ties on the order key exist and you want a
row-precise window, state \texttt{ROWS BETWEEN} explicitly.

\paragraph{Gaps and islands.} The signal is any mention of consecutive values,
streaks, or runs. An \emph{island} is a maximal run of consecutive values. The
trick: subtract a dense row number from the value. Across a consecutive run
both sides increase by one per step, so the difference stays constant, and it
jumps exactly at a gap. Group by the difference to collapse each run.

\begin{lstlisting}[style=sql, caption={Islands of consecutive ids}]
WITH numbered AS (
    -- DISTINCT first so duplicate ids do not desync the row number
    SELECT DISTINCT log_id,
           ROW_NUMBER() OVER (ORDER BY log_id) AS rn
    FROM logs
)
SELECT MIN(log_id) AS island_start,
       MAX(log_id) AS island_end,
       COUNT(*)    AS length
FROM numbered
GROUP BY log_id - rn
ORDER BY island_start;
\end{lstlisting}

The \texttt{DISTINCT} is load-bearing. If \texttt{log\_id} repeats, \texttt{rn}
still advances on the duplicate while the value does not, so
\texttt{log\_id - rn} shifts and splits an island that should stay whole. For
consecutive \emph{dates} rather than integers, offset by an interval:
\texttt{log\_date - rn * INTERVAL `1 day'} is constant within a run.

\begin{lstlisting}[style=sql, caption={Gaps via LEAD}]
SELECT log_id + 1  AS missing_from,
       next_id - 1 AS missing_to
FROM (
    SELECT log_id,
           LEAD(log_id) OVER (ORDER BY log_id) AS next_id
    FROM logs
) t
WHERE next_id - log_id > 1;   -- values missing between two present ids
\end{lstlisting}

\subsubsection{Deduplication}

\paragraph{DISTINCT.} \texttt{DISTINCT} compares the entire select list as a
tuple. \texttt{SELECT DISTINCT a, b} returns every distinct \emph{pair}, so a
value in \texttt{a} can still repeat as long as the paired \texttt{b} differs.
It is a whole-row operator with no way to say ``dedupe on this column and keep
the rest of the row.'' It runs after \texttt{SELECT} and before
\texttt{ORDER BY}, which is why every \texttt{ORDER BY} expression must also
appear in the select list, and it treats two \texttt{NULL}s as duplicates.
\texttt{COUNT(DISTINCT col)} is a separate mechanism that dedupes the input to
one aggregate and has no effect on the rows returned.

\newpage

\paragraph{DISTINCT ON.} A PostgreSQL extension that keeps the first row of
each group sharing the same value of the parenthesized expressions, returning
all selected columns of that surviving row. ``First'' is defined by
\texttt{ORDER BY}, whose leading terms must match the \texttt{DISTINCT ON}
expressions. Omit the \texttt{ORDER BY} and the surviving row is arbitrary.

\begin{lstlisting}[style=sql, caption={Most recent close per ticker}]
SELECT DISTINCT ON (ticker)
       ticker, trade_date, close
FROM prices
ORDER BY ticker, trade_date DESC;   -- leading term must be ticker
\end{lstlisting}

\vspace{-1ex}

This is the compact form of the rank-then-filter pattern from the window
section. The \texttt{ROW\_NUMBER} equivalent needs a subquery and an outer
\texttt{WHERE rn = 1}. Which to reach for:

\begin{itemize}[label=$\circ$]
  \item \textbf{Top-1 per group in Postgres:} \texttt{DISTINCT ON} is shorter
  and generally at least as fast, since it stops at the first row of each
  group.
  \item \textbf{Top-$N$ for $N > 1$, or portability:} \texttt{ROW\_NUMBER}
  extends to \texttt{rn <= 3} and runs on any engine. \texttt{DISTINCT ON} is
  Postgres-only.
  \item \textbf{Ties:} \texttt{DISTINCT ON} always emits exactly one row per
  group, settled by the remaining \texttt{ORDER BY} terms or arbitrarily. To
  keep every tied row, use \texttt{RANK() = 1}.
\end{itemize}

Because the \texttt{ORDER BY} is dictated by the deduplication key, presenting
the result in a different order means wrapping the query in an outer
\texttt{SELECT ... ORDER BY}.

\subsubsection{String Manipulation}

Positions are 1-based, and every function returns a new value rather than
mutating in place. The recurring uses are normalizing a dirty join key, parsing
a delimited field, and folding a group of strings into one.

\paragraph{Concatenation.} The operator \texttt{a || b} propagates
\texttt{NULL}: if either side is \texttt{NULL} the whole expression is
\texttt{NULL}. \texttt{CONCAT(a, b, ...)} treats \texttt{NULL} as an empty
string. \texttt{CONCAT\_WS(sep, ...)} takes the separator first, joins the rest
with it, and skips \texttt{NULL} arguments without leaving a doubled separator.
Choosing \texttt{||} for user-facing text is a common way to silently blank out
a whole column.

\paragraph{Case and whitespace.} \texttt{UPPER} and \texttt{LOWER} are the
standard folds, \texttt{INITCAP} capitalizes the first letter of each word and
lowercases the rest. \texttt{TRIM} strips characters from the ends only, spaces
by default, with \texttt{LEADING}, \texttt{TRAILING}, or \texttt{BOTH} to pick
a side. Interior whitespace survives, so collapsing double spaces needs
\texttt{REGEXP\_REPLACE(s, '\textbackslash s+', ' ', 'g')}.
\texttt{LPAD(s, n, fill)} and \texttt{RPAD} pad to a fixed width and truncate
when the input is already longer than \texttt{n}.

\vspace{2ex}

\begin{lstlisting}[style=sql, caption={Normalizing a join key}]
SELECT
    LOWER(TRIM(email))                      AS email_key,
    SPLIT_PART(LOWER(TRIM(email)), '@', 2)  AS domain,
    LPAD(account_id::text, 8, '0')          AS padded_id,
    CONCAT_WS(', ', last_name, first_name)  AS sort_name
FROM users;
\end{lstlisting}

\newpage

\paragraph{Extraction and search.} \texttt{SUBSTRING(s FROM a FOR b)} (or
\texttt{SUBSTR(s, a, b)}) takes \texttt{b} characters starting at position
\texttt{a}. \texttt{LEFT(s, n)} and \texttt{RIGHT(s, n)} take from the ends,
and a negative \texttt{n} drops that many characters from the opposite end.
\texttt{STRPOS(s, sub)} returns the 1-based index of the first occurrence and
\texttt{0} when the substring is absent, so a found/not-found test compares
against \texttt{0}. \texttt{SPLIT\_PART(s, delim, n)} pulls the \texttt{n}th
field of a delimited string and returns an empty string when that field does
not exist, worth wrapping in \texttt{NULLIF(..., '')}.

\begin{lstlisting}[style=sql, caption={Parsing a delimited name field}]
SELECT full_name,
       SPLIT_PART(full_name, ' ', 1)              AS first_name,
       NULLIF(SPLIT_PART(full_name, ' ', 2), '')  AS last_name,
       LENGTH(full_name)
         - LENGTH(REPLACE(full_name, ' ', ''))    AS space_count
FROM people;
\end{lstlisting}

\vspace{-1ex}

The last expression is the occurrence-counting idiom: remove every copy of the
target and measure how much the string shrank. For a needle longer than one
character, divide the difference by \texttt{LENGTH(needle)}.

\paragraph{Replacement.} \texttt{REPLACE(s, from, to)} substitutes every
literal occurrence. \texttt{TRANSLATE(s, from, to)} maps character to character
by position and deletes any character in \texttt{from} that has no counterpart
in \texttt{to}, a fast way to strip a known character set.
\texttt{REGEXP\_REPLACE(s, pattern, repl, flags)} handles the general case and
replaces only the \emph{first} match unless the \texttt{`g'} flag is passed, a
frequent source of half-cleaned output.

\paragraph{Pattern matching.} \texttt{LIKE} uses \texttt{\%} for any run of
characters and \texttt{\_} for exactly one, and it is case sensitive.
\texttt{ILIKE} is the case-insensitive form. For real expressiveness the POSIX
operators apply: \texttt{\textasciitilde} matches, \texttt{\textasciitilde*}
matches case-insensitively, and \texttt{!\textasciitilde} negates. Anchors
\texttt{\^{}} and \texttt{\$} matter, since an unanchored pattern matches
anywhere in the string.

\begin{lstlisting}[style=sql, caption={Regex cleaning and validation}]
SELECT phone,
       REGEXP_REPLACE(phone, '[^0-9]', '', 'g')  AS digits,
       phone ~ '^\+?[0-9]{10,15}$'               AS looks_valid,
       email ILIKE '%@gmail.com'                 AS is_gmail
FROM contacts;
\end{lstlisting}

\vspace{-1ex}

\paragraph{Splitting and aggregating.} \texttt{STRING\_AGG(expr, delim ORDER BY
...)} collapses a group into one delimited string and is the string analogue of
\texttt{SUM}. The \texttt{ORDER BY} lives inside the call, so the output order
is deterministic. The inverse direction expands one row into many:
\texttt{UNNEST(STRING\_TO\_ARRAY(s, delim))} or
\texttt{REGEXP\_SPLIT\_TO\_TABLE(s, pattern)} in the \texttt{FROM} clause.

\begin{lstlisting}[style=sql, caption={Collapse a group, then expand it back}]
SELECT department,
       STRING_AGG(name, ', ' ORDER BY name) AS members
FROM employees
GROUP BY department;
-- one row per tag stored in a comma-separated column
SELECT p.id, TRIM(tag) AS tag
FROM products p,
     UNNEST(STRING_TO_ARRAY(p.tags, ',')) AS tag;
\end{lstlisting}


\newpage

\subsubsection{Command Reference}

Signatures are PostgreSQL. Optional arguments are shown in brackets.

\begin{table}[H]
\centering
\footnotesize
\renewcommand{\arraystretch}{1.4}
\rowcolors{2}{gray!10}{white}
\begin{tabularx}{\textwidth}{>{\raggedright\arraybackslash}p{0.20\textwidth} >{\raggedright\arraybackslash}p{0.225\textwidth} >{\raggedright\arraybackslash}X}
\toprule
\textit{Command} & \textit{Arguments} & \textit{What it does} \\
\midrule
\textbf{WHERE} & boolean predicate & Filters individual rows before grouping. Cannot see aggregates, window functions, or select aliases. \\
\textbf{GROUP BY} & column or expression list & Collapses rows to one output row per distinct key combination. \\
\textbf{HAVING} & boolean predicate & Filters whole groups after aggregation. The only legal home for an aggregate predicate. \\
\textbf{DISTINCT} & applies to the select list & Removes duplicate rows compared across the entire select list. Treats two \texttt{NULL}s as equal. \\
\textbf{DISTINCT ON} & \texttt{(expr, ...)} & Postgres. Keeps one row per distinct value of the expressions, chosen by the leading \texttt{ORDER BY} terms, which must match. \\
\textbf{ORDER BY} & \texttt{expr [ASC|DESC] [NULLS FIRST|LAST]} & Sorts the output. Postgres places \texttt{NULL}s last under \texttt{ASC} and first under \texttt{DESC}. \\
\textbf{LIMIT / OFFSET} & \texttt{n} & Caps the number of rows returned, or skips the first \texttt{n}. \\
\textbf{WITH} & \texttt{name AS (query)} & Common table expression. Names a subquery for reuse. \texttt{WITH RECURSIVE} walks hierarchies. \\
\textbf{CASE} & \texttt{WHEN cond THEN val [ELSE val] END} & Row-level branching. Returns the first match, or \texttt{NULL} when nothing matches and \texttt{ELSE} is absent. \\
\textbf{FILTER} & \texttt{(WHERE cond)} after an aggregate & Restricts which rows the aggregate consumes. Postgres shorthand for \texttt{SUM(CASE ...)}. \\
\textbf{EXISTS / NOT EXISTS} & subquery & Semi-join and anti-join. Null-safe, unlike \texttt{IN} and \texttt{NOT IN}. \\
\bottomrule
\end{tabularx}
\caption*{Query clauses and structure}
\end{table}

\vspace{-2ex}

\begin{table}[H]
\centering
\footnotesize
\renewcommand{\arraystretch}{1.4}
\rowcolors{2}{gray!10}{white}
\begin{tabularx}{\textwidth}{>{\raggedright\arraybackslash}p{0.20\textwidth} >{\raggedright\arraybackslash}p{0.225\textwidth} >{\raggedright\arraybackslash}X}
\toprule
\textit{Command} & \textit{Arguments} & \textit{What it does} \\
\midrule
\textbf{COUNT(*)} & none & Counts every row, including all-\texttt{NULL} ones. \\
\textbf{COUNT(col)} & one column & Counts rows where \texttt{col} is non-\texttt{NULL}. \\
\textbf{COUNT(\allowbreak DISTINCT col)} & one column & Counts distinct non-\texttt{NULL} values. \\
\textbf{SUM / AVG /\allowbreak MIN / MAX} & one expression & Skip \texttt{NULL}s. \texttt{AVG} divides by the non-\texttt{NULL} count. \\
\textbf{STRING\_AGG} & \texttt{(expr, delim [ORDER BY ...])} & Concatenates a group into one delimited string. \\
\textbf{ARRAY\_AGG} & \texttt{(expr [ORDER BY ...])} & Collects group values into an array. \\
\textbf{COALESCE} & any number of values & Returns the first non-\texttt{NULL} argument, left to right. \\
\textbf{NULLIF} & two values & Returns \texttt{NULL} when the arguments are equal, otherwise the first. Guards division by zero. \\
\textbf{GREATEST / LEAST} & list of values & Largest or smallest value across columns within one row. Ignores \texttt{NULL}s. \\
\textbf{IS [NOT] DISTINCT FROM} & two values & Null-safe comparison that treats two \texttt{NULL}s as equal. \\
\textbf{IS [NOT] NULL} & one value & The only valid null test. \texttt{= NULL} is always unknown. \\
\bottomrule
\end{tabularx}
\caption*{Aggregates and \texttt{NULL} handling}
\end{table}

\vspace{-2ex}

\begin{table}[H]
\centering
\footnotesize
\renewcommand{\arraystretch}{1.4}
\rowcolors{2}{gray!10}{white}
\begin{tabularx}{\textwidth}{>{\raggedright\arraybackslash}p{0.20\textwidth} >{\raggedright\arraybackslash}p{0.225\textwidth} >{\raggedright\arraybackslash}X}
\toprule
\textit{Command} & \textit{Arguments} & \textit{What it does} \\
\midrule
\textbf{OVER} & \texttt{(PARTITION BY ... ORDER BY ... [frame])} & Turns a function into a window function evaluated per row without collapsing the result set. \\
\textbf{ROW\_NUMBER()} & none & Sequential number within the partition. Ties broken arbitrarily. \\
\textbf{RANK()} & none & Competition rank. Ties share a rank and leave gaps afterward. \\
\textbf{DENSE\_RANK()} & none & Ties share a rank with no gaps. \\
\textbf{NTILE(n)} & bucket count & Splits the partition into \texttt{n} roughly equal buckets and labels each row. \\
\textbf{LAG(col[, k, default])} & \texttt{k} defaults to 1 & Value \texttt{k} rows earlier in the ordered partition. \texttt{NULL} past the partition edge unless a default is supplied. \\
\textbf{LEAD(col[, k, default])} & same as \texttt{LAG} & Value \texttt{k} rows ahead. \\
\textbf{FIRST\_VALUE /\allowbreak LAST\_VALUE} & \texttt{col} & First or last value in the current frame. \texttt{LAST\_VALUE} needs an explicit frame to see the full partition. \\
\textbf{NTH\_VALUE(col, n)} & column, position & The \texttt{n}th value in the frame. \\
\textbf{SUM / AVG / COUNT ... OVER} & aggregate plus window spec & Running or rolling aggregate computed per row. \\
\textbf{ROWS BETWEEN} & \texttt{a PRECEDING} to \texttt{b FOLLOWING}, \texttt{CURRENT ROW}, \texttt{UNBOUNDED} & Physical frame counted in rows. Use it for fixed moving windows. \\
\textbf{RANGE BETWEEN} & same bound keywords & Value-based frame that includes every row tied on the \texttt{ORDER BY} key. This is the default frame. \\
\bottomrule
\end{tabularx}
\caption*{Window functions and frames}
\end{table}

\newpage

\begin{table}[H]
\centering
\footnotesize
\renewcommand{\arraystretch}{1.4}
\rowcolors{2}{gray!10}{white}
\begin{tabularx}{\textwidth}{>{\raggedright\arraybackslash}p{0.20\textwidth} >{\raggedright\arraybackslash}p{0.225\textwidth} >{\raggedright\arraybackslash}X}
\toprule
\textit{Command} & \textit{Arguments} & \textit{What it does} \\
\midrule
\textbf{\texttt{||}} & two strings & Concatenation. Yields \texttt{NULL} if either side is \texttt{NULL}. \\
\textbf{CONCAT} & any number of values & Concatenation treating \texttt{NULL} as an empty string. \\
\textbf{CONCAT\_WS} & separator, then values & Joins arguments with a separator and skips \texttt{NULL}s. \\
\textbf{LENGTH / CHAR\_LENGTH} & one string & Character count. \texttt{OCTET\_LENGTH} gives bytes. \\
\textbf{UPPER / LOWER /\allowbreak INITCAP} & one string & Case folding. \texttt{INITCAP} capitalizes each word and lowercases the rest. \\
\textbf{TRIM} & \texttt{[LEADING|TRAILING|BOTH] [chars] FROM s} & Strips characters from the ends only, spaces by default. \texttt{LTRIM}, \texttt{RTRIM}, \texttt{BTRIM} are shorthands. \\
\textbf{LPAD / RPAD} & \texttt{(s, n [, fill])} & Pads to width \texttt{n}, truncating when \texttt{s} is already longer. \\
\textbf{LEFT / RIGHT} & \texttt{(s, n)} & First or last \texttt{n} characters. A negative \texttt{n} drops that many from the opposite end. \\
\textbf{SUBSTRING} & \texttt{(s FROM a FOR b)} or \texttt{SUBSTR(s, a, b)} & Slice of length \texttt{b} starting at 1-based position \texttt{a}. \\
\textbf{STRPOS / POSITION} & \texttt{(s, sub)} / \texttt{(sub IN s)} & 1-based index of the first occurrence, \texttt{0} when absent. \\
\textbf{REPLACE} & \texttt{(s, from, to)} & Replaces every literal occurrence. \\
\textbf{TRANSLATE} & \texttt{(s, from\_set, to\_set)} & Maps characters by position and deletes those with no counterpart. \\
\textbf{SPLIT\_PART} & \texttt{(s, delim, n)} & The \texttt{n}th delimited field. Returns an empty string when the field is missing. \\
\textbf{STRING\_TO\_\allowbreak ARRAY / UNNEST} & \texttt{(s, delim)} / \texttt{(array)} & Splits into an array, then expands an array into rows. \\
\textbf{REVERSE / REPEAT} & \texttt{(s)} / \texttt{(s, n)} & Reverses a string, or repeats it \texttt{n} times. \\
\textbf{LIKE / ILIKE} & pattern with \texttt{\%} and \texttt{\_} & Wildcard match, case sensitive and case insensitive respectively. \\
\textbf{\textasciitilde{} \textasciitilde* !\textasciitilde{}} & POSIX regex & Regex match, case-insensitive match, and negated match. \\
\textbf{TO\_CHAR} & \texttt{(value, format)} & Formats a number or timestamp as text, e.g.\ \texttt{'YYYY-MM-DD'}. \\
\textbf{CAST / \texttt{::}} & \texttt{(expr AS type)} / \texttt{expr::type} & Type conversion. \texttt{::numeric} also fixes integer division. \\
\bottomrule
\end{tabularx}
\caption*{String functions}
\end{table}

